\section{Neurocognitive methods: the primary score boundary}
\label{sec:p19-ncm-annex}
\begin{annexitem}{annex:p19.ncm-score}{thm:p19-ncm-score}{Vogel's two score contexts and their exact domains}
\begin{proof}
Suppose \(F\ne T\). At \(D=T\), the numerator equals
the denominator, so \(\alpha=1\); at \(D=F\), the
numerator vanishes, so \(\alpha=0\). If \(F>T\),
the inequalities \(0\leq(F-D)/(F-T)\leq1\) are
equivalent to \(0\leq F-D\leq F-T\), hence
\(T\leq D\leq F\). If \(F<T\), multiplication
by the negative denominator reverses both inequalities and yields
\(F\leq D\leq T\). These prove necessity and sufficiency
for either ordering. For \(F=T\), division by zero is not
defined, even when the numerator also vanishes. No value is
assigned to that missing physiological contrast.

In the distinct capacity context, \(H, F\in[0, 1]\) imply
\(-1\leq H-F\leq1\); multiplication by \(S>0\)
gives the displayed range. Furthermore \(K\geq0\) is
equivalent to \(H-F\geq0\), hence \(H\geq F\).
The endpoints are attainable at \((H, F)=(0, 1), (1, 0)\).
These are algebraic properties of the formulas. They neither
establish the latent capacity model nor identify which brain
region generates an amplitude or a behavioral response.
\end{proof}
\end{annexitem}
